%\documentclass[a4paper,12pt]{article}  % standard LaTeX, 12 point type
\documentclass[12pt, a4paper]{article}

\usepackage{algpseudocode}
\usepackage{algorithm}
\usepackage{algorithmicx}

\usepackage{geometry}
\usepackage{amsfonts,latexsym}
\usepackage{amsthm}
\usepackage{amssymb}
\usepackage[utf8]{inputenc} % Кодировка
\usepackage[english,russian]{babel} % Многоязычность
\usepackage{mathtools}
\usepackage{hyperref}
\usepackage{tikz}
\usepackage{dsfont}
\usepackage{multicol}
\usetikzlibrary{fit,calc,automata,positioning}

\theoremstyle{definition}
\newtheorem{definition}{Определение}[section]
\newtheorem{example}{Пример}[section]
\newtheorem{theorem}{Теорема}[section]
\newtheorem{proposition}[theorem]{Proposition}
\newtheorem{lemma}[theorem]{Лемма}
\newtheorem{corollary}[theorem]{Corollary}
\newtheorem{conjecture}[theorem]{Conjecture}


% unnumbered environments:

\theoremstyle{remark}
\newtheorem*{remark}{Remark}
\newtheorem*{notation}{Notation}
\newtheorem*{note}{Note}



\setlength{\parskip}{5pt plus 2pt minus 1pt}
%\setlength{\parindent}{0pt}


\algtext*{EndWhile}% Remove "end while" text
\algtext*{EndIf}% Remove "end if" text
\algtext*{EndFor}% Remove "end for" text
\algtext*{EndFunction}% Remove "end function" text


\usepackage{color}
\usepackage{listings}
\usepackage{caption}
\usepackage{graphicx}
\usepackage{ucs}

\graphicspath{{pics/}}

\geometry{left=2cm}
\geometry{right=1.5cm}
\geometry{top=2cm}
\geometry{bottom=2cm}




%\lstnewenvironment{algorithm}[1][]
%{
%    \lstset{
%        frame=tB,
%        numbers=left,
%        mathescape=true,
%        numberstyle=\small,
%        basicstyle=\small,
%        inputencoding=utf8,
%        extendedchars=\true,
%        keywordstyle=\color{black}\bfseries,
%        keywords={,function, procedure, return, datatype, function, in, if, else, for, foreach, while, denote, do, and, then, assert,}
%        numbers=left,
%        xleftmargin=.04\textwidth,
%        #1 % this is to add specific settings to an usage of this environment (for instnce, the caption and referable label)
%    }
%}
%{}

\newcommand{\tab}[1][0.3cm]{\ensuremath{\hspace*{#1}}}

\newcommand{\rvline}{\hspace*{-\arraycolsep}\vline\hspace*{-\arraycolsep}}

\newcommand{\derives}[1][*]{\xRightarrow[]{#1}}
\newcommand{\first}[1][1]{\textsc{first}_{#1}}
\newcommand{\follow}[1][1]{\textsc{follow}_{#1}}

\setcounter{MaxMatrixCols}{20}


\tikzset{
%->, % makes the edges directed
%>=stealth’, % makes the arrow heads bold
node distance=4cm, % specifies the minimum distance between two nodes. Change if necessary.
%every state/.style={thick, fill=gray!10}, % sets the properties for each ’state’ node
initial text=$ $, % sets the text that appears on the start arrow
}

\tikzstyle{symbol_node} = [shape=rectangle, rounded corners, draw, align=center]
\tikzstyle{prod_node} = [shape=rectangle, draw, align=center]

\tikzset{
    between/.style args={#1 and #2}{
         at = ($(#1)!0.5!(#2)$)
    }
}

%every node/.style = {shape=rectangle, rounded corners,
%      draw, align=center,
%      top color=white, bottom color=blue!20}

\title{Практика программироавния. Заметки.}
\author{Семён Григорьев}
\date{\today}

\begin{document}
\maketitle
\newpage
\tableofcontents
\newpage

%\section{План лекций}
\begin{enumerate}
  \item Введение. Базовые определения. Обзор курса.
  \item Регулярные языки, конечные автоматы (детерминированные, недетерминированные), регулярные выражения. Детерминизация, $\varepsilon-$замыкание, минимизация.
  \item Взаимные преобразования способов задания.
  \item Теоретико-языковые свойства регулярных языков. Лемма о накачке, замкнутость относительно операций.
  \item Алгоритмы вычисления операций.
  \item Грамматики, переписывающие системы. КС-граммтики (обыкновенные граммтики). Вывод в граммтике, неоднозначные грамматики, существенно неоднозначные языки, дерево вывода.
  \item Лево(право)-линейные грамматики и регулярные языки.
  \item КС граммтики и КС языки. Лемма о накачке, замкнутотсть относительно операций, проверка пустоты.
  \item !!!
  \item !!!
  \item !!!
  \item !!!
  \item !!!
  \item !!!
  \item !!!
  \item !!!
\end{enumerate}


\section{Лекция 1: Разреженные матрицы}

\subsection{Форматы разреженных матриц}

Важно то, что их дофига. И разные хороши для разных задач.

Рассмотрм матрицу 

$$
M=
\begin{pmatrix}
1 & 0 & 0 & 0 & 3 & 0 & 9 & 0 \\
0 & 0 & 2 & 0 & 0 & -1 & 0 & 0 \\
0 & 0 & 0 & 7 & 0 & 0 & 0 & 0 \\
0 & 6 & 0 & 0 & 0 & 0 & 0 & 5 \\
0 & 3 & 4 & 0 & 0 & 0 & 0 & 0 \\
0 & 0 & 0 & 0 & 2 & 0 & 8 & 0 \\
0 & 0 & -3 & 0 & 0 & 0 & 0 & -7 \\
0 & 0 & 0 & 0 & -4 & 0 & 0 & 0 
\end{pmatrix}.
$$
В ней всего 15 ненулевыз элементов (NNZ(M) = 15).
Попробуем представить её в разных форматах.

\textbf{Покоординатный (COO).} Часто он же \textit{triple store}. Просто коллекция троек вида \textit(строка, столбец, значение) только для ненулевых элементов.
\begin{align*}
M_{\textit{COO}} = [&(0,0,1);(0,4,3);(0,6,9);(1,2,2);(1,5,-1);(2,3,7);(3,1,6);(3,7,5);(4,1,3);\\
           &(4,2,4);(5,4,2);(5,6,8); (6,2,-3);(2,7,-7);(7,4,-4) ]
\end{align*}

Потратили $3*NNZ(M)$ памяти.

Формат простой. Можно предварительно сортировать по столбцам или строкам. Достаточно просто удалять и добавлять элементы. Не самый оптимальный по доступу и по дополнительной памяти.

\textbf{Compressed row storage (CRS).} Заведём 3 массива: \textit{val} для хранения значений, \textit{col\_id} для хранения номера столбца, \textit{row\_ptr} для указателя на начало строки в \textit{col\_id}.

\begin{center}
$$
M_{\textit{CRS}} = (\textit{val},\textit{col\_id},\textit{row\_ptr} )
$$
\begin{tabular}{|c||c|c|c|c|c|c|c|c|c|c|c|c|c|c|c|}
\hline
val     & 1 & 3 & 9 & 2 & -1 & 7 & 6 & 5 & 3 & 4 & 2 & 8 & -3 & -7 & -4\\
col\_id & 0 & 4 & 6 & 2 & 5  & 3 & 1 & 7 & 1 & 2 & 4 & 6 & 2  & 7  & 4\\
\hline
\end{tabular}

\begin{tabular}{|c||c|c|c|c|c|c|c|c|}
\hline
row\_ptr & 0 & 3 & 5 & 6 & 8  & 10 & 12 & 14 \\
\hline
\end{tabular}
\end{center}

Потратили $2*NNZ(M) + n$ памяти. Сложнее построить, сложно удалять и добавлять. Достаточно просто последовательно обходить, лучше по памяти и по обращению к элементу. Наиболее часто встречается в современных библиотеках.

\textbf{Compressed column storage (CCS)} и \textbf{compressed dioganal storage (CDS).} Вариации на тему CRS. Второй --- для диоганальных матриц.

\textbf{Блочный CRS (CCS, CDS).} Когда мы знаем, что матрица имеет блочную структуру, а каждый блок --- разреженная матрица. Пример --- результат тензорного произведения.

\textbf{Quad-tree.}

Предположим, что $n=2^k$. Если это не так, то выравниваем нулями. После этого строим рекурсивно дерево: дели матрицу на 4 блока, корень --- это текущая матрица, листья --- 4 полученных блока. Если в блоке все элементы нули, то сразу рисуем лист типа \textit{(этот блок --- ноль)}, иначе продолжаем разбиение. И так до тех пор пока не получатся блоки из одного элемента.


$$
M=
\left[
\begin{array}{cccc|cccc}
1 & 0 & 0 & 0 & 3 & 0 & 9 & 0 \\
0 & 0 & 2 & 0 & 0 & -1 & 0 & 0 \\
0 & 0 & 0 & 7 & 0 & 0 & 0 & 0 \\
0 & 6 & 0 & 0 & 0 & 0 & 0 & 5 \\
\hline
0 & 3 & 4 & 0 & 0 & 0 & 0 & 0 \\
0 & 0 & 0 & 0 & 2 & 0 & 8 & 0 \\
0 & 0 & -3 & 0 & 0 & 0 & 0 & -7 \\
0 & 0 & 0 & 0 & -4 & 0 & 0 & 0 
\end{array}
\right]
$$

$$
M_{(0,0)}=
\left[
\begin{array}{cc|cc}
1 & 0 & 0 & 0 \\
0 & 0 & 2 & 0 \\
\hline
0 & 0 & 0 & 7 \\
0 & 6 & 0 & 0 \\
\end{array}
\right],
M_{(0,1)}=
\left[
\begin{array}{cc|cc}
3 & 0 & 9 & 0 \\
0 & -1 & 0 & 0 \\
\hline
0 & 0 & 0 & 0 \\
0 & 0 & 0 & 5 \\ 
\end{array}
\right]
$$

$$
M_{(1,0)}=
\left[
\begin{array}{cc|cc}
0 & 3 & 4  & 0 \\
0 & 0 & 0  & 0 \\
\hline
0 & 0 & -3 & 0 \\
0 & 0 & 0  & 0  
\end{array}
\right],
M_{(1,1)}=
\left[
\begin{array}{cc|cc}
 0  & 0 & 0 & 0 \\
 2  & 0 & 8 & 0 \\
 \hline
 0  & 0 & 0 & -7 \\
 -4 & 0 & 0 & 0 
\end{array}
\right]
$$

$$
M_{(0,0)(0,0)}=
\left[
\begin{array}{c|c}
1 & 0 \\
\hline
0 & 0 
\end{array}
\right],
M_{(0,0)(0,1)}=
\left[
\begin{array}{c|c}
0 & 0 \\
\hline
2 & 0 
\end{array}
\right],
M_{(0,0)(1,0)}=
\left[
\begin{array}{c|c}
0 & 0 \\
\hline
0 & 6 
\end{array}
\right],
M_{(0,0)(1,1)}=
\left[
\begin{array}{c|c}
0 & 7 \\
\hline
0 & 0 
\end{array}
\right]
$$

$$
M_{(0,1)(0,0)}=
\left[
\begin{array}{c|c}
3 & 0  \\
\hline
0 & -1 \\
\end{array}
\right],
M_{(0,1)(0,1)}=
\left[
\begin{array}{c|c}
9 & 0 \\
\hline
0 & 0 \\
\end{array}
\right],
M_{(0,1)(1,1)}=
\left[
\begin{array}{c|c}
0 & 0 \\
\hline
0 & 5 \\ 
\end{array}
\right]
$$

$$
M_{(1,0)(0,0)}=
\left[
\begin{array}{c|c}
0 & 3 \\
\hline
0 & 0 
\end{array}
\right],
M_{(1,0)(0,1)}=
\left[
\begin{array}{c|c}
4  & 0 \\
\hline
0  & 0 
\end{array}
\right],
M_{(1,0)(1,1)}=
\left[
\begin{array}{c|c}
 -3 & 0 \\
 \hline
 0  & 0  
\end{array}
\right]
$$

$$
M_{(1,1)(0,0)}=
\left[
\begin{array}{c|c}
 0  & 0 \\
 \hline
 2  & 0 
\end{array}
\right],
M_{(1,1)(0,1)}=
\left[
\begin{array}{c|c}
0 & 0 \\
\hline
8 & 0 
\end{array}
\right],
M_{(1,1)(1,0)}=
\left[
\begin{array}{c|c}
 0  & 0\\
 \hline
 -4 & 0
\end{array}
\right],
M_{(1,1)(1,1)}=
\left[
\begin{array}{c|c}
 0  & -7 \\
 \hline
 0 & 0 
\end{array}
\right]
$$

\resizebox{.99\textwidth}{!}{
\begin{tikzpicture}[sibling distance=27em,
  every node/.style = {shape=rectangle, rounded corners,
    draw, align=center,
    top color=white, bottom color=blue!20}]]
  \node {$M$}
    child {[sibling distance=6em] 
        node {$M_{(0,0)}$} 
        child { [sibling distance=1.5em] 
            node {$M_{(0,0)(0,0)}$} 
            child {node {$1$}}
            child {node {$0$}}
            child {node {$0$}}
            child {node {$0$}} }
        child { [sibling distance=1.5em]
            node {$M_{(0,0)(0,1)}$} 
            child {node {$0$}}
            child {node {$0$}}
            child {node {$2$}}
            child {node {$0$}} }
        child { [sibling distance=1.5em]
            node {$M_{(0,0)(1,0)}$}
            child {node {$0$}}
            child {node {$0$}}
            child {node {$0$}}
            child {node {$6$}} }
        child { [sibling distance=1.5em]
            node {$M_{(0,0)(1,1)}$}
            child {node {$0$}}
            child {node {$7$}}
            child {node {$0$}}
            child {node {$0$}} } }
    child { [sibling distance=7em] 
        node {$M_{(0,1)}$}
        child { [sibling distance=2em] 
            node {$M_{(0,1)(0,0)}$}
            child {node {$3$}}
            child {node {$0$}}
            child {node {$0$}}
            child {node {$-1$}} }
        child { [sibling distance=1.5em] 
            node {$M_{(0,1)(0,1)}$}
            child {node {$9$}}
            child {node {$0$}}
            child {node {$0$}}
            child {node {$0$}} }
        child { node {$0$} }
        child { [sibling distance=1.5em] 
            node {$M_{(0,1)(1,1)}$}
            child {node {$0$}}
            child {node {$0$}}
            child {node {$0$}}
            child {node {$5$}} } }
    child { [sibling distance=6em]
        node {$M_{(1,0)}$} 
        child { [sibling distance=1.5em] 
            node {$M_{(1,0)(0,0)}$}
            child {node {$0$}}
            child {node {$3$}}
            child {node {$0$}}
            child {node {$0$}} }
        child { [sibling distance=1.5em] 
            node {$M_{(1,0)(0,1)}$}
            child {node {$4$}}
            child {node {$0$}}
            child {node {$0$}}
            child {node {$0$}} }
        child { node {$0$} }
        child { [sibling distance=2em] 
            node {$M_{(1,0)(1,1)}$} 
            child {node {$-3$}}
            child {node {$0$}}
            child {node {$0$}}
            child {node {$0$}} } }
    child { [sibling distance=7.5em]
        node {$M_{(1,1)}$} 
        child { [sibling distance=1.5em] 
            node {$M_{(1,1)(0,0)}$}
            child {node {$0$}}
            child {node {$0$}}
            child {node {$2$}}
            child {node {$0$}} }
        child { [sibling distance=1.5em] 
            node {$M_{(1,1)(0,1)}$}
            child {node {$0$}}
            child {node {$0$}}
            child {node {$8$}}
            child {node {$0$}} }
        child { [sibling distance=2em] 
            node {$M_{(1,1)(1,0)}$}
            child {node {$0$}}
            child {node {$0$}}
            child {node {$-4$}}
            child {node {$0$}} }
        child { [sibling distance=2em] 
            node {$M_{(1,1)(1,1)}$}
            child {node {$0$}}
            child {node {$-7$}}
            child {node {$0$}}
            child {node {$0$}} } };
\end{tikzpicture}
}

Относительно сложно построить, легко разделить на части (хорошо для распределённых вычислений), достаточно быстрый доступ к элементу, простой рекурсивный алгоритм перемножения.

Часто используется в композиции с CRS или другими форматами: до какого-то момента строим дередо, а потом в листьях храним сравнительно большие подматрицы в CRS.

\subsection{Умножение разреженных матриц}

Временная сложность оценивается относительно количества ненулевых элементов (NNZ) либо во входных матрицах, либо в выходной матрице, либо и там и там.

При использовании GPGPU возникают проблемы с необходимостью динамически выделять память. Поэтому используют два шага: оценка NNZ  в результирующей матрице с последующим выделением нужного количества памяти, собственно перемножение. 
\section{Практика 1}

Детали о том, как будет проходить практика.

\subsection{Григорьев С.В.}

Немного про описания языков.
Пописать языковые уравнения, граммтики.
Посмотреть на операции над языками.

Постановка задачи на весь семестр.

Запросы к графовым базам данных.
Контекст задачи, примеры графовых БД (RedisGraph, Neo4j, ...), задача о путях впринципе.

Ссылка на второй конспект.

Задача: реализовать свою "графовую миниБД".

Реализация: оформление, инструменты, языки.

\begin{itemize}
  \item Ограничений на язык реализации нет.
  \item Ограничений на использование библиотек нет. Главное --- не нарушать лицензии и чтобы можно было вносить изменения в библиотеку (при необходимости).
  \item Каждый создёт под решение репозиторий на GitHub и снабжает его всем необходимым: readme, лицензия, CI-сборка с тестированием, инструкции по локальному развёртыванию.
  \item Разработка ведётся в отдельной ветке и когда очередная часть задачи готова к сдаче --- делаем pull request в master и добавляем меня (gsvgit) в ревьюверы.
\end{itemize}


Задачи на дом.
\begin{enumerate}
  \item Выбрать язык программирования, на котором будет вестись разработка.
  \item Создать репозиторий на GitHub.
  \item Настроить CI-сборку и тестирование.
  \item Реализовать подгрузку графов из RDF используя готовые библиотеки.
\end{enumerate}


\section{Спектральная теория графов. Введение}

Матрица смежности.

Матрица Кирхгофа. Оператор Лапласа. Для неориентированного графа без кратных рёбер и петель.
\begin{definition}
Пусть неориентированный граф без кратных рёбер и петель (простой граф) $G=\langle V, E \rangle, |V| = n$. Тогда матрица Кирхгофа $ K=(k_{i,j})_{n \times n}$. 
$$ k_{i,j}:={\begin{cases}\deg(v_{i})&{\text{при}}\ i=j,\\-1&{\text{при}}\ (v_{i},v_{j})\in E(G),\\0&{\text{в противном случае}}.\end{cases}}
$$
\end{definition}

$K = D - A$, где $A$ --- матрица смежности графа, а $D$ --- матрица, на диагонали которой строят степени вершин, а остальные элементы равны нулю.


\begin{example}[Пример графа и его матрицы Кирхгофа]
Пусть дан граф:
  \begin{center}
  \begin{tikzpicture}[on grid, auto]
     \node[state] (q_0)   {$0$};
     \node[state] (q_1) [above right=1.4cm and 1.0cm of q_0] {$1$};
     \node[state] (q_2) [right=2.0cm of q_0] {$2$};
     \node[state] (q_3) [right=2.0cm of q_2] {$3$};
      \path[-]
      (q_0) edge (q_1)
      (q_1) edge (q_2)
      (q_2) edge (q_0)
      (q_2) edge (q_3);
  \end{tikzpicture}
  \end{center}
  $$ D =
  \left({
  \begin{array}{rrrrrr}
  2 & 0 & 0 & 0 \\
  0 & 2 & 0 & 0 \\
  0 & 0 & 3 & 0 \\
  0 & 0 & 0 & 1 \\
  \end{array}
  }\right)
$$
$$ A =
  \left({
  \begin{array}{rrrrrr}
  0 & 1 & 1 & 0 \\
  1 & 0 & 1 & 0 \\
  1 & 1 & 0 & 1 \\
  0 & 0 & 1 & 0 \\
  \end{array}
  }\right)
$$

$$ K = D - A =
  \left({
  \begin{array}{rrrrrr}
  2  & -1 & -1 & 0  \\
  -1 & 2  & -1 & 0  \\
  -1 & -1 & 3  & -1 \\
  0  & 0  & -1 & 1  \\
  \end{array}
  }\right)
$$
\end{example}

Определитель матрицы.


\begin{definition}[Дополнительный минор]
$M_{i,j}$ --- дополнительный минор, определитель матрицы, получающейся из исходной матрицы $A$ путём вычёркивания $i$-й строки и $j$-го столбца.
\end{definition}


\begin{definition}[Определитель матрицы $2 \times 2$]
Для матрицы $2\times 2$ определитель вычисляется как:

$$\Delta ={\begin{vmatrix}a&c\\b&d\end{vmatrix}} = ad - bc$$

\end{definition}

\begin{definition}[Определитель матрицы $N \times N$]

$$\Delta =\sum _{j=0}^{n-1}(-1)^{j}a_{0,j}{M}_{0,j}$$
, где $M_{0,j}$ --- дополнительный минор к элементу $a_{0,j}$. 


\end{definition}

Это было разложение по строке и, вообще говоря, подобная операция может быть проделяна для любой строки.
Аналогично можно использовать разложение по столбцу.


\begin{definition}[Определитель матрицы $3 \times 3$]
\begin{align*}
&\Delta =
{\begin{vmatrix}
  a_{0,0}&a_{0,1}&a_{0,2}\\
  a_{1,0}&a_{1,1}&a_{1,2}\\
  a_{2,0}&a_{2,1}&a_{2,2}
 \end{vmatrix}}
 =
  a_{0,1}{\begin{vmatrix}a_{1,1}&a_{1,2}\\a_{2,1}&a_{2,2}\end{vmatrix}}
 -a_{0,2}{\begin{vmatrix}a_{1,0}&a_{1,2}\\a_{2,0}&a_{2,2}\end{vmatrix}}
 +a_{0,3}{\begin{vmatrix}a_{1,0}&a_{1,1}\\a_{2,0}&a_{2,1}\end{vmatrix}} 
 =\\
&a_{0,0}a_{1,1}a_{2,2}-a_{0,0}a_{1,2}a_{2,1}-a_{0,1}a_{1,0}a_{2,2}+a_{0,1}a_{1,2}a_{2,0}+a_{0,2}a_{1,0}a_{2,1}-a_{0,2}a_{1,1}a_{2,0}
\end{align*}
\end{definition}


\begin{definition}[Алгебраическое дополнение]
Алгебраическим дополнением элемента $ a_{i,j}$ матрицы $A$ называется число $A_{i,j}=(-1)^{i+j}M_{i,j}$,
где $M_{i,j}$ --- дополнительный минор.
\end{definition}

\begin{example}[Определитель]
Найдём определитель матрицы Кирхгофа для нашего графа. Будем использовать разложение по 3-й строке.
\begin{align*}
&\Delta\left({
  \begin{array}{rrrrrr}
  2  & -1 & -1 & 0  \\
  -1 & 2  & -1 & 0  \\
  -1 & -1 & 3  & -1 \\
  0  & 0  & -1 & 1  \\
  \end{array}
  }\right) = (-1)^5 (-1) {\begin{vmatrix}2&-1&0\\-1&2&0\\-1&-1&-1\end{vmatrix}} + (-1)^6 1 {\begin{vmatrix}2&-1&-1\\-1&2&-1\\-1&-1&3\end{vmatrix}}=\\
  &1((2\cdot2\cdot-1) - (2\cdot0\cdot-1) - (-1\cdot-1\cdot-1) + (-1\cdot0\cdot-1) + (0\cdot-1\cdot-1) - (0\cdot2\cdot-1)) +\\
  & 1((2\cdot2\cdot3) - (2\cdot-1\cdot-1) - (-1\cdot-1\cdot3) + (-1\cdot-1\cdot-1) + (-1\cdot-1\cdot-1) - (-1\cdot2\cdot-1) ) =\\
  & (-4 + 1) + (12 - 2 - 3 - 1 - 1 - 2) = -3 + 3 = 0 
\end{align*}
\end{example}


\begin{theorem}[Матричная теорема об остовных деревьях]
Пусть $G$ --- связный простой граф с матрицей Кирхгофа $K$. Все алгебраические дополнения матрицы Кирхгофа $K$ равны между собой и их общее значение равно количеству остовных деревьев графа $G$.
\end{theorem}

\begin{example}[Количество остовных деревьев]
Из примера выше, значения миноров равно 3.

Деревья:
\begin{center}
  \begin{tikzpicture}[on grid, auto]
     \node[state] (q_0)   {$0$};
     \node[state] (q_1) [above right=1.4cm and 1.0cm of q_0] {$1$};
     \node[state] (q_2) [right=2.0cm of q_0] {$2$};
     \node[state] (q_3) [right=2.0cm of q_2] {$3$};
      \path[-]
      (q_0) edge (q_1)
      (q_2) edge (q_0)
      (q_2) edge (q_3);
  \end{tikzpicture}
  \end{center}

  \begin{center}
  \begin{tikzpicture}[on grid, auto]
     \node[state] (q_0)   {$0$};
     \node[state] (q_1) [above right=1.4cm and 1.0cm of q_0] {$1$};
     \node[state] (q_2) [right=2.0cm of q_0] {$2$};
     \node[state] (q_3) [right=2.0cm of q_2] {$3$};
      \path[-]
      (q_1) edge (q_2)
      (q_2) edge (q_0)
      (q_2) edge (q_3);
  \end{tikzpicture}
  \end{center}

  \begin{center}
  \begin{tikzpicture}[on grid, auto]
     \node[state] (q_0)   {$0$};
     \node[state] (q_1) [above right=1.4cm and 1.0cm of q_0] {$1$};
     \node[state] (q_2) [right=2.0cm of q_0] {$2$};
     \node[state] (q_3) [right=2.0cm of q_2] {$3$};
      \path[-]
      (q_0) edge (q_1)
      (q_1) edge (q_2)
      (q_2) edge (q_3);
  \end{tikzpicture}
  \end{center}
\end{example}


Сумма элементов каждой строки (столбца) матрицы Кирхгофа равна нулю: $\sum _{i=0}^{|V|-1}k_{i,j}=0$.

Определитель матрицы Кирхгофа равен нулю: $\Delta K = 0$.

Собственные числа и собственные вектора.

Нам понядобится поле $F$. 

\begin{definition}[Собственный вектор]
Ненулевой вектор $x$ называется собственным вектором матрицы $A$ для некоторого элемента $\lambda \in F$, если $Ax = \lambda x$
\end{definition}

\begin{definition}[Собственное число (Собственное значение)]
Собственным числом матрицы $A$ называется такое $\lambda \in F$, что существует ненулевое решение уравнения $Ax = \lambda x$.
\end{definition}

Как видно, собственные числа и собственные вектора ``ходят парами''.

\begin{example}[Собственные числа и вектора]

$$
A = \left({
  \begin{array}{rrrrrr}
  2  & -1 & -1 & 0  \\
  -1 & 2  & -1 & 0  \\
  -1 & -1 & 3  & -1 \\
  0  & 0  & -1 & 1  \\
  \end{array}
  }\right)
$$

По определению $Ax = \lambda x$. 
$$Ax - \lambda x = 0$$
$$(A - \lambda E) x = 0$$
$$
(\left({
  \begin{array}{rrrrrr}
  2  & -1 & -1 & 0  \\
  -1 & 2  & -1 & 0  \\
  -1 & -1 & 3  & -1 \\
  0  & 0  & -1 & 1  \\
  \end{array}
  }\right) -  \left({
  \begin{array}{rrrrrr}
  \lambda  & 0 & 0 & 0  \\
  0 & \lambda  & 0 & 0  \\
  0 & 0 & \lambda  & 0 \\
  0  & 0  & 0 & \lambda  \\
  \end{array}
  }\right)) x = 0
$$

Данное уравнение имеет ненулевое решение тогда и только тогда, когда $|A - \lambda E| = 0$
\begin{align*}
|A - \lambda E| =
  \begin{vmatrix}
  2 - \lambda & -1 & -1 & 0  \\
  -1 & 2 - \lambda & -1 & 0  \\
  -1 & -1 & 3 - \lambda & -1 \\
  0  & 0  & -1 & 1 - \lambda \\
  \end{vmatrix} = \lambda^4-8\lambda^3+19\lambda^2-12\lambda
\end{align*}
То есть надо решить уравнение 
$$
\lambda^4-8\lambda^3+19\lambda^2-12\lambda = 0
$$

$$
\lambda^4-8\lambda^3+19\lambda^2-12\lambda  = \lambda (\lambda^3 - 8\lambda^2 + 19\lambda - 12) = \lambda(\lambda - 1)(\lambda^2 - 7\lambda + 12) = \lambda(\lambda - 1)(\lambda - 3)(\lambda - 4) = 0
$$

Корни: $\lambda \in \{0,1,3,4\}$.

Далее для каждого совственного числа нужно найти соответсвующий вектор. Для этого решаем системы линейных уравнений (метод Гаусса в помощь).
$$(A - 0\cdot E) x = 0$$
$$(A - 1\cdot E) x = 0$$
$$(A - 3\cdot E) x = 0$$
$$(A - 4\cdot E) x = 0$$

$$x_0 = \left(\begin{array}{r}1\\1\\1\\1\end{array}\right)$$
$$x_1 = \left(\begin{array}{r}-\frac{1}{2}\\-\frac{1}{2}\\0\\1\end{array}\right)$$
$$x_2 = \left(\begin{array}{r}-1\\1\\0\\0\end{array}\right)$$
$$x_3 = \left(\begin{array}{r}1\\1\\-3\\1\end{array}\right)$$


\end{example}

\begin{definition}[Спектр графа]
Спектром графа называется упорядоченное по возростанию мультимножество собственных значений его матрицы смежности.
\end{definition}

Так как мы говорим о неориентированном графе, то собственные значения всегд вещественные числа (почему?).

Хотя матрица смежности и зависит от нумерации вешин, спектр является инвариантом графа (почему?).

Следствие: изоморфные графы имеют одинаковый спектр.

Графы с одинаковым спектром --- изоспектральные (или коспектральные).

\begin{theorem}
Изоморфные графы всегда изоспектральны. Обратное не верно (изоспектральные графы не обязательно изоморфны).
\end{theorem}

\subsection{Задачи}
\begin{enumerate}
\item Доказать формулу Кэли, пользуясь матричной теоремой об остовных деревьях. Формула Кэли даёт оценку числа остовных деревьев полного графа $K_{n}$: $n^{n-2}$.
\item Доказать, пользуясь матричной теоремой об остовных деревьях, что число остовных деревьев полного двудольнoгo графа $K_{m,n}$ равно $m^{n-1}\cdot n^{m-1}$.
\item Доказать, что спектр является инвариантом графа.
\end{enumerate}
\section{Практика 2}

\subsection{Григорьев С.В.}

Построение ДКА по регулярному варажению.
Парсер регулярных выражений. Немного про AST и прочее.

Пересечение ДКА и НКА (без $\varepsilon$-переходов): библиотека и руками через произведение Кронекера (прямое произведение). Сравнение производительности.

Представление результата в виде графа и регулярного выражения.








\section{Лекция 3}

	Ещё раз произменения: в реквесте должно быть только то, что непосредственно относится к сдаваемой домашке.

	Ещё раз про функции, про то, как выделять и разделять функциональность, не надо запихивать всё в одну функцию. Про то, где должны быть проверки.

	Про консоль.

	Про обработку крайних случаев. Про исключения.

	Про тесты и ошибки: нашёл ошибку --- создал тест.

	Про стиль кодирования: про пробелы вокруг скобок и операций, про отступы и переводы строк. Про соглашения о наименовании. camlCase CamlCase


	Про единицы измерения.

    Базовые структуры данных, алгоритмы и их выражение в F\#. Функция. Рекурсия и итерация.  Базовые типы и основы работы с ними: матрицы, массивы, списки, структуры. 


    Числа Фибоначчи.
\section{Практика 3}

\subsection{Григорьев С.В.}

Пересечение автоматов --- это тензорное произведение матриц смежности. Пример.

Про коммутативность пересечения и некоммутативность тензорного произведения.

Домашнее задание.
\begin{enumerate}
	\item Реализовать консольный клиент, позволяющий
	\begin{enumerate}
		\item загрузить RDF-файл
		\item вывести список меток рёбер
		\item задать к загруженному графу регулярный запрос с возможностью указать представление результата: пустота ответа, автомат сдампить в файл в формате DOT (\url{https://www.graphviz.org/doc/info/lang.html}), пара (кол-во рёбер, кол-во вершин) в результирующем автомате
	    \item выйти из клиента.
    \end{enumerate}
	\item Подгрузку RDF и выполнение запросов реализовать на основе уже существующей функциональности.
	\item Провести замеры производительности на графах из репозитория \url{https://github.com/JetBrains-Research/CFPQ_Data}. Графы брать из подпапки \verb|data/graphs/RDF|. Так как в графах присутствуют одинаковые отношения, то можно один и тот же запрос выполнять на всех графах. Отчёт оформить в виде раздела в  README репозитория в виде таблицы. 

	\textit{Эти эксперименты проводятся локально! Не надо таскать репозиторий с графами за собой. Для тестов клиента использовать маленькие синтетические RDF.}
	\item Реализовать необходимые тесты на работоспособность клиента.
\end{enumerate} 


\section{Лекция 4}
    Рассказать про то, что выделенные значения --- это плохо. Немного про исключительные ситуации.
    Про создание проектов. Про версии пакетов и вообще версии артефактов. Про то, что заимствование кода не поощрается. Тем более неправомерное заимствование. Про классный пример форматирования \verb|x*x*x + x*x + x + 1|.

    Работа с файлами.
    
    Сортировки: пузырьком, вставкой, Хоара. Различные сценарии использования: поддержание отсортированного набора, сортировка всего набора целиком. Некоторые особенности реализации: наивная функциональная реализация Хоара, реализация на массиве.
    
    Основы машинного представления данных. Представления чисел. Представление чисел с плавающей точкой. Проблемы переполнения. Битовые операции. Строки, кодировки.

\section{Домашняя работа 4}

Во всех задачах на сортировку необходимо реализовать чтение массива из файла и печать результата в файл. Функции чтения и записи необходимо переиспользовать.

В задачах на битовые операции продолжаем рабоать с консолью: чтение данных с консоли и печать результата туда же.

Для данной домашней работы необходимо создать отдельный проект.

При создании тестов необходимо, в задачах на сортировку, убедиться, что, во-первых, сортировки ведут себя одинаоково на одинаковых данных, во-вторых, что они ведут себя так же, как системные сортировки для соответствующих коллекций. Для задачи о запаковке и распаковке надо проверить, что реализованные функции являются взаимно обратными. FsCheck (testProperty) в помощь.

\begin{enumerate}
    \item \textbf{(1 балл)} Реализовать сортировку пузырьком массива. 
    \item \textbf{(1 балл)} Реализовать сортировку пузырьком списка.
    \item \textbf{(1 балл)} Реализовать быструю сортировку для списка.
    \item \textbf{(1 балл)} Реализовать быструю сортировку для массива.
    \item \textbf{(1 балл)} Реализовать запаковку двух 32-битных чисел в одно 64-битное и распаковку обратно. 
    \item \textbf{(1 балл)} Реализовать запаковку четырёх 16-битных чисел в одно 64-битное и распаковку обратно. 
\end{enumerate}


\section{Лекция 5}
 
\subsection{Нормальня форма Хомского (НФХ)}

\begin{definition}
КС грамматика находится в нормальной форме Хомского если любое правило имеет один из трёх видов:
\begin{enumerate}
    \item $S \to \varepsilon$
    \item $N_i \to t_j$
    \item $N_i \to N_j N_k$, $N_j \neq S, N_k \neq S$
\end{enumerate}
\end{definition}
Важно: стартовый нетерминал не встречается в правых частях правил, $\varepsilon$-продукция только для стартового нетерминала.

\begin{note}
Любую КС грамматику можно преобразвать к нормальной форме Хомского.
\end{note}

Преобразование в НФХ. Шаги.
\begin{enumerate}
    \item Устранение длинных правил.
    \item Устранение $\varepsilon$-правил.
    \item Устранение цепных правил.
    \item Устранение бесполезных нетерминалов
    \begin{enumerate}
        \item Удаление непорождающих нетерминалов
        \item Удаление недостижимых нетерминалов
    \end{enumerate}
    \item Устранение продукций с правой частью длины 2, содержащей терминалы. 
\end{enumerate}

Надо не забыть добавить новый стартовый нетерминал, если нужно: чтобы вывести из него $\varepsilon$ и чтобы не встречался в правых частях правил. 

\textbf{Важно.}
Порядок применения шагов преобразования важен.
\begin{enumerate}
    \item Второй шаг можно поднять наверх, но это приведёт к более существенному разростанию результирующей граммтики.
    \item Подшаги шага 4 нельзя менять местами. Попробуйте поприменять их к граммтике:

\begin{align*}
S \to & A B \mid a \\
A \to & b
\end{align*}
    
\end{enumerate}

\href{https://neerc.ifmo.ru/wiki/index.php?title=%D0%9D%D0%BE%D1%80%D0%BC%D0%B0%D0%BB%D1%8C%D0%BD%D0%B0%D1%8F_%D1%84%D0%BE%D1%80%D0%BC%D0%B0_%D0%A5%D0%BE%D0%BC%D1%81%D0%BA%D0%BE%D0%B3%D0%BE}{Материалы по преобразованию в НФХ.}


\subsection{Лемма о накачке для КС языков}

\begin{theorem}
Пусть $L$ --- контекстно-свободный язык над алфавитом $\Sigma$, тогда существует такое $n$, что для любого слова $\omega \in L$, $|\omega| \geq n$ найдутся слова $u,v,x,y,z\in \Sigma^*$, для которых верно: $uvxyz = \omega, vy\neq \varepsilon,|vxy|\leq n$ и для любого $k \geq 0$  $uv^kxy^kz \in L$.
\end{theorem}

Идея доказательства леммы о накачке.

\begin{enumerate}
    \item Для любого КС языка можно найти грамматику в нормальной форме Хомского.
    \item Очевидно, что если брать достаточно длинные цепочки, то в дереве вывода этих цепочек, на пути от корня к какому-то листу обязательно будет нетерминал, встречающийся минимум два раза. Если $m$ --- количество нетерминалов в НФХ, то длины $2^{m+1}$ должно хватить. Это и будет $n$ из леммы.
    \item Возьмём путь, на котором есть хотя бы дважды повторяется некоторый нетерминал. Скажем, это нетерминал  $N_1$. Пойдём от листа по этому пути. Найдём первое появление $N_1$. Цепочка, задаваемая поддеревом для этого узла --- это $x$ из леммы.
    \item Пойдём дальше и найдём второе появление $N_1$. Цепочка, задаваемая поддеревом для этого узла --- это $vxy$ из леммы.
    \item Теперь мы можем копировать кусок дерева между этими повторениями $N_1$ и таким образом накачивать исходную цепочку.
\end{enumerate}

Надо только проверить выполение ограничений на длины.

\begin{figure}
\centering
\includegraphics[width=0.5\textwidth]{figures/pumping_tree_1.pdf}
\caption{Разбиение цепочки для леммы о накачке}
\label{fig:pumping1}
\end{figure}

\begin{figure}
\centering
\includegraphics[width=0.5\textwidth]{figures/pumping_tree_2.pdf}
\caption{Пример накачки цепочки с рисунка~\ref{fig:pumping1}}
\end{figure}


\href{https://neerc.ifmo.ru/wiki/index.php?title=%D0%9B%D0%B5%D0%BC%D0%BC%D0%B0_%D0%BE_%D1%80%D0%B0%D0%B7%D1%80%D0%B0%D1%81%D1%82%D0%B0%D0%BD%D0%B8%D0%B8_%D0%B4%D0%BB%D1%8F_%D0%9A%D0%A1-%D0%B3%D1%80%D0%B0%D0%BC%D0%BC%D0%B0%D1%82%D0%B8%D0%BA}{Материалы по лемме о накачке для КС языков.}

Проверить неконтекстно-свободность языка $L=\{a^nb^nc^n \mid n>0\}$.

\section{Практика 5}

\subsection{Григорьев С.В.}

Преобразование в нормальную форму Хомского.

Формат входа:
\begin{enumerate}
\item Одна продукция на строку.
\item Продукция --- это список терминалов и нетерминалов через пробел, начинающийся с нереминала (левая часть продукции).
\item Нетерминалы --- заглавные буквы с опциональным числовым суффиксом.
\item Терминалы --- строчные буквы с опциональным числовым суффиксом.
\item Специальный символ \verb|eps| для обобзначения $\varepsilon$.
\end{enumerate}

Пример входя, описывающего граммтику $S \to a S b S \mid \varepsilon$:

\begin{verbatim}
S a S b S
S eps 
\end{verbatim}

Домашнее задание.
\begin{enumerate}
    \item Реализовать преобразование в нормальную форму Хомского. На входе файл с граммтикой, на выходе --- файл с граммтикой в НФХ в том же формате, что и вход.
\end{enumerate} 



\section{Лекция 6}
 
\subsection{Свойства замкнутости КС языков}

Глава 2.6 \href{https://github.com/YaccConstructor/articles/blob/master/InProgress/Formal_langs_CFPQ_course_notes/Formal_lang_CFPQ_course_notes.pdf}{конспекта}.

\subsection{Алгоритм CYK}

Глава 4.1 ``Алгоритм CYK'' \href{https://github.com/YaccConstructor/articles/blob/master/InProgress/Formal_langs_CFPQ_course_notes/Formal_lang_CFPQ_course_notes.pdf}{конспекта}.
%b\section{Практика 5}

\subsection{Григорьев С.В.}

Преобразование в нормальную форму Хомского.

Формат входа:
\begin{enumerate}
\item Одна продукция на строку.
\item Продукция --- это список терминалов и нетерминалов через пробел, начинающийся с нереминала (левая часть продукции).
\item Нетерминалы --- заглавные буквы с опциональным числовым суффиксом.
\item Терминалы --- строчные буквы с опциональным числовым суффиксом.
\item Специальный символ \verb|eps| для обобзначения $\varepsilon$.
\end{enumerate}

Пример входя, описывающего граммтику $S \to a S b S \mid \varepsilon$:

\begin{verbatim}
S a S b S
S eps 
\end{verbatim}

Домашнее задание.
\begin{enumerate}
    \item Реализовать преобразование в нормальную форму Хомского. На входе файл с граммтикой, на выходе --- файл с граммтикой в НФХ в том же формате, что и вход.
\end{enumerate} 


\section{Лекция 7}

Пример рефакторинга разбора аргументов командной строки. И то же самое с тестами.

 
Понятие типа данных. Системы типов: статические, динамические, строгие, нестрогие. Примеры языков с разными системами типов. Понятие о разной выразительности (``мощности'') систем.
Пример кода на \href{http://www.fstar-lang.org/tutorial/}{$F^\star$}.

Приведение типов: автоматическое, ручное. %Вывод типов по Хиндли-Милнеру.

Обобщённые типы данных. Понятие о полиморфизме.
    
Алгебраические типы данных: кортежи, Discriminated Unions (рамеченные объединения). Примеры на F\#. Единицы измерения.

\section{Домашняя работа 6}

\begin{enumerate}
    \item \textbf{(1 балл)} Предположим, что мы храним булевы матрицы в виде списка координат ячеек, значение которых \texttt{true}. Необходимо реализовать соответствующие типы: единицы измерения для строк и столбцов, пара ``строка-столбец'', список пар ``строка-столбец''. 
    \item \textbf{(2 балла)} Реализовать функцию, перемножающую две матрицы, заданных в формате, описанном в предыдущей задаче. Не забыть проверку корректности входных данных. Реализовать подгрузку матриц из файла и запись результата в файл. Файлы с данными и результатом указываются через консоль. Формат хранения матрицы из $m$ строк и $n$ столбцов: в файле $m$ строк, каждая строка состоит из $n$ символов 0 или 1.
\end{enumerate} 



\section{Лекция 8}
 
Нисходящий синтаксический анализ: рекурсивный спуск и LL(k).


Глава 7, 8.1 и 8.2 \href{https://github.com/YaccConstructor/articles/blob/master/InProgress/Formal_langs_CFPQ_course_notes/Formal_lang_CFPQ_course_notes.pdf}{конспекта}.

\section{Практика 7}

\subsection{Григорьев С.В.}

Алгоритмы на основе линейной алгебры.

Для реализации предлагается использовать следующие библиотеки. Так как с булевыми не везде хорошо, то будем использовать те типы, которые поддерживаются: \verb|Int|, \verb|float| и т.д.
\begin{itemize}
    \item Для языка Python --- \href{https://docs.scipy.org/doc/scipy/reference/sparse.html}{разреженные матрицы в scipy} и соответствующие опреации работы с ними: \href{https://docs.scipy.org/doc/scipy-0.14.0/reference/generated/scipy.sparse.kron.html}{\textbf{scipy.sparse.kron}} и \href{https://docs.scipy.org/doc/scipy/reference/generated/scipy.sparse.csr_matrix.html}{обычное матричное произведение}. Предпочтительный формат разреженных матриц --- CSR. 
    \item Для языка Kotlin ---\href{http://la4j.org/}{la4j}. Операции: \href{http://la4j.org/apidocs/org/la4j/operation/ooplace/OoPlaceKroneckerProduct.html}{кронекер} и  \href{http://la4j.org/apidocs/org/la4j/operation/ooplace/OoPlaceMatricesMultiplication.html}{обычное умножение}. 
\end{itemize} 

Поэлементное сложение есть и там и там.

Для алгоритма на матричном умножении всё точно так же, как и в предыдущей ДЗ для Хеллингса. 

Для тензорного произведения расширим формат представления входной грамматики.
Одна строка на нетерминал. Терминалы, нетерминалы, $\varepsilon$ обозначаются как и раньше.
Как и раньше, левая часть от правой отделена пробелом. В правой части можно использовать конструкции регулярных выражений: альтернатива, звезда клини, групперцющие скобки. Этот набор можно расширять по своему усмотрению.


Пример входа, описывающего граммтику $S \to (a S b)* \mid \varepsilon$:

\begin{verbatim}
S (a S b)* | eps 
\end{verbatim}


Домашнее задание. Время на выполнение --- две недели. Один из алгоритмов --- на первую, оставшийся и эксперименты --- на вторую.
\begin{enumerate}
    \item Реализовать алгоритм, основанный на матричном умножении. На вход принимаются два файла: с граммтикой и входным графом. В результирующий файл печатается граммтика в ослабленной НФХ (с которой непосредсвенно работал алгоритм) и множество пар достижимых вершин для стартового нетерминала (одна пара на строку, две вершины через пробел).
    \item Реализовать алгоритм, основанный на тензорном произведении. На вход принимается файл с граммтикой и файл с графом. В результирующий файл печатается матрица смежности рекурсивного автомата (с которым непосредсвенно работал алгоритм, построчно, элементы разделены пробелом, пустая ячейка обозначается символом '.') и множество пар достижимых вершин для стартового нетерминала (одна пара на строку, две вершины через пробел).
    \item Сравнить производительность трёх реализованных алгоритмов (Хеллингс, матричное произведение, тензорное произведение). Результат --- описание эксперимента и таблица сравнения в readme.  
\end{enumerate} 



\section{Лекция 9}
 
Нисходящий синтаксический анализ: GLL и его применение для поиска путей с КС ограничениями.


Глава 8.3 и 8.4 \href{https://github.com/YaccConstructor/articles/blob/master/InProgress/Formal_langs_CFPQ_course_notes/Formal_lang_CFPQ_course_notes.pdf}{конспекта}.

\section{Домашняя работа 8}

При создании новых структур данных необходимо расширять библиотеку, созданную в предыдущей домашней работе.

\begin{enumerate}

    \item \textbf{(5 балла)} Используя тип MyList из предыдущей домашней работы, реализовать целочисленную длинную арифметику: операции сложения, умножения, вычитания, целочисленного деления. 

    \item \textbf{(2 балл)} Реализовать представление разреженных матриц в виде дерева квадрантов. Реализовать функцию поэлементного сложения двух матриц в таком формате.
    
    \item \textbf{(3 балл)} Реализовать функцию умножения двух матриц в формате дерева квадрантов.
    
    \item \textbf{(3 балл)} Реализовать функцию тензорного умножения двух матриц в формате дерева квадрантов.
    
    \item \textbf{(5 балла)} Реализовать построение транзитивного замыкания ориентированного графа через произведение матриц. Использовать представление матриц из первой задачи. Визуализировать результат с помощью GraphViz: исходный граф и выделенные рёбра, появившиеся в результате транзитивного замыкания. Для задания графа использовать формат из задачи 2 6-й домашней работы.
    
    \item \textbf{(5 балла)} Реализовать вычисление кратчайших путей между всеми парами вершин в ориентированном графе. Использовать представление матриц из первой задачи. Визуализировать результат с помощью GraphViz: исходный граф и выделенные рёбра со значением кратчайшего пути между соответствующей парой вершин. Для задания графа использовать формат, аналогичный формату из задачи 2 6-й домашней работы.
\end{enumerate}

\section{Лекция 10}
 
Восходящий синтаксический анализ: алгоритмы семейства LR.


Глава 9.1 \href{https://github.com/YaccConstructor/articles/blob/master/InProgress/Formal_langs_CFPQ_course_notes/Formal_lang_CFPQ_course_notes.pdf}{конспекта}.


\section{Лекция 11}

Проделжение про FsLex+FsYacc, полноценный пример интерпретатора. 
 

Интерпретация и компиляция. Особенности, разновидности интерпретаторов, основные шаги.  Особенности, разновидности компиляторов, основные шаги. Примеры, пример реализации простого интерпретатора.

Устройство сред разработки и компиляторов, интерпретаторов: общие шаги, классические возможности, JIT/AOT. Примеры из .NET, F\#, JVM, LLVM.

F\# $\to$ IL $\to$ C\# $\to$ ASM

\begin{verbatim}
module M

type Lst<'a> =
    | Nill
    | Cons of 'a * Lst<'a>
    
let rec map f l =
    match l with
    | Nill -> Nill
    | Cons (x,tl) -> Cons(f x, map f tl)

Cons(1,Nill)
|> map ((+)1) 
|> printfn "%A"

\end{verbatim}
Этот код на \href{https://sharplab.io/#v2:DYLgZgzgNALiCWwoBMQGoA+BbA9sgrsAKYAEAsgLABQ1MAngA6kAyEMAPAOQCGAfCQF5qJESQwkAcomDDR4gMI4AdhBI4wJHiQBUJVhx69ZI6sRgkATkQDGJLNwYkNwQcbvcY1gBYkXAd3gYLzdxKWAXAFp+MJkqUTESRRUSAAoAD1hgAEoSKMTlCBSNDPdHDRhs6mokwoBGKBis6gx+e0cUlLQs2pzm/gYLeCUYMCUSACIAUgBBcbdqIA==}{sharplab.io}


\begin{verbatim}
module M

let f x = 
    let h y = if y / 2 = 0 then y * x else y + x
    if x > 5 then x + 2 else h (x * 6)
    
let rec g y = f (y + g 2)

\end{verbatim}
Этот код на \href{https://sharplab.io/#v2:DYLgZgzgNAJiDUAfAtgexgV2AUwAQFkBYAKBJwBdcxcAPXAXlxNxdwtwAtcBPB3AS2q8A9LgBMfAAy5yHbADseuAFS1c2YBDy94tZq0FqAfLgCsMuYrq6JGrZ1wAKOqoBsASn0sy2SgCdsAGNcAHMlRmpHHVDxT2IgA=}{sharplab.io}
\section{Практика 9}

\subsection{Григорьев С.В.}

Абстрактный и конкретный синтаксис языка запросов к графам, семантика языка запросов к гарфам.

Абстрактный синтаксис скрывает детали синтаксического анализа и конкретного ``текстового'' представления синтаксиса, оставляя только важную для дальнейшего анализа иформацию о структуре программы. 

Первая, минималистичная, версия языка запросов к графам.

\begin{verbatim}
type script: Seq of List<stmt>

type stmt:
  | Connect of String
  | Select of obj_expr * graph_expr

type graph_expr =
  | Intersect of graph_expr * graph_expr 
  | Query of pattern
  | GraphName of String

type obj_expr:
  | Edges   
  | Count of obj_expr
  
type pattern: 
  | smb of String
  | star of pattern
  | plus of pattern
  | alt of pattern * pattern
  | seq of pattern * pattern
  | option of pattern

\end{verbatim}

\newcommand{\sem}[1]{\llbracket #1 \rrbracket}

Тперь поговорим о семантике данного языка. 

Семантика --- отображение из программ в семантический домен.
Программы обычно представляются как абстрактные синтаксические деревья.
Семантический домен --- то, что описывает смысл программы в контексте решаемой задачи.

Семантика бывает разной.
Классические примеры: динамическая семантика (вычисление программ) и статическая семантика (вывод типов).

Доменом динамической семантики арифметических выражений является множество функций из функции, сопоставляющей переменным целочисленные значения, в целые числа.

Семантика часто обозначается скобками $\llbracket$ и $\rrbracket$.

Семаника арифметических выражений:

$$\sem{\cdot} : (X \rightarrow \mathbb{N}) \rightarrow \mathbb{N}$$


\begin{align*}
  \sem{n}(\sigma) &= n, n \in \mathbb{N} \\
  \sem{v}(\sigma) &= \sigma(v) \\
  \sem{e_1 \oplus e_2}(\sigma) &= \sem{e_1}(\sigma) \otimes \sem{e_2}(\sigma) \\
  & \\
  \sem{-e}(\sigma) &= \sem{0-e}(\sigma)
\end{align*}

Здесь $\oplus$ --- конструкторы бинарных операций, $\otimes$ --- соответствующие арифметические функции.

В нашем случае, исполнение скрипта на получившимся языке достаточно естественно описывается состоянием из текущей базы данных, текущего графа, и текущего результата вычислений. Значение этой тройки и будут семантическим доменом нашего языка.

При обработке непустой последовательности инструкций сначала вычисляем результат выполнения первой инструкции, затем на нем --- всех остальных.

$$
\infer[]{c \xrightarrow{Seq \ (hd::tl)} c''}{c \xrightarrow{hd} c', \ \ c' \xrightarrow{Seq \ tl} c''}
$$

При обработке команды \verb|Connect| сохраняем её аргумент в конфигурации как путь к текущей базе данных.
$$
\infer[]{\langle db, graph, res\rangle \xrightarrow{Connect \ db'} \langle db', graph, res\rangle}{}
$$

При обработке инструкции \verb|Select| сперва вычисляется граф, из которого будут выбираться объекты, затем вычисляются сами объекты.

$$
\infer[]{c \xrightarrow{Select \ obj \ graph} c''}{c \xrightarrow{graph} c', \ \ c' \xrightarrow{obj} c''}
$$

При обработке инструкции \verb|Graph| загружаем в окружение граф из файла с именем, соотвествующим аргументу этой инструкции, лежащем в базе (директории), соответствующей значению \verb|cur_db|.

$$
\infer[]{\langle db, graph, res \rangle \xrightarrow{\textit{GraphName} \ name} \langle db ,loadGraph(String.concat(db,name)), res\rangle}{}
$$

При обработке инструкции \verb|Intersect| сперва вычисляем оба подвыражения, затем пересекаем получившиеся графы как конечные автоматы. Считаем все вершины стартовыми и финальными одновременно.

$$
\infer[]{\langle db, graph, res\rangle \xrightarrow{\textit{Intersect} \ graph_1 \ graph_2} \langle db, \textit{intersectFA}(graph', graph''), res\rangle}{\langle db, graph,res\rangle \xrightarrow{graph_1} \langle db, graph', res \rangle, \ \ \langle db, graph,res \rangle \xrightarrow{graph_2} \langle db, graph'', res \rangle}
$$


При обработке инструкции \verb|Query| необходимо построить минимальный детерминированный автомат по регулярному выражению, задаваемому шаблоном.

$$
\infer[]{\langle db, graph, res \rangle \xrightarrow{\textit{Query} \ pattern} \langle db, \textit{buildMinDFA}(pattern), res \rangle}{}
$$

При обработке инструкции \verb|Edges| в текущий результат вычислений записывается множество троек, соотвествующих рёбрам текущего графа.

$$
\infer[]{\langle db, graph, res \rangle \xrightarrow{\textit{Edges}} \langle db, graph, \{(v_i,e,v_j) \mid (v_i,e,v_j) \in graph.Edges\}\rangle}{}
$$

При обработке инструкции \verb|Count| сперва вычисляется подвыражене, а затем в текущий результат вычислений записывается размер результата вычисления подвыражения, если результат был множеством.

$$
\infer[type\_of(res') \equiv Set ]{\langle db, graph, res \rangle \xrightarrow{\textit{Count} \ expr} \langle db, graph, Set.count(res') \rangle}{\langle db, graph, res \rangle \xrightarrow{expr} \langle db, graph, res' \rangle}
$$

--------------------------------------------------------------------------------




Расширенная версия.

\begin{verbatim}
type script = Seq of List<stmt>

type stmt =
  | Connect of String
  | NamedPattern of String * pattern
  | Select of obj_expr * graph_expr

type graph_expr =
  | Intersect of graph_expr * graph_expr 
  | Query of pattern
  | GraphName of String
  | SetStartAndFinal of vertices * vertices * graph_expr

type vertices =
  | Set of Set<int>
  | Range of int * int
  | None   

type obj_expr =
  | Edges
  | Filter of cond * obj_expr
  | Count of obj_expr
  
type cond = 
  | Cond of String * String * String * bool_expr

type bool_expr = 
  | LblIs   of String * String
  | IsStart of String
  | IsFinal of String

type pattern = 
  | Term of String
  | Nonterm of String
  | Star of pattern
  | Plus of pattern
  | Alt of pattern * pattern
  | Seq of List<pattern>
  | Option of pattern

\end{verbatim}


Для описания семантики для добавленных конструкций потребуется изменить домен --- в него необходимо добавить grammar (текущую граммтику).


При обработке инструкции \verb|NamedPattern| добавляем в \textit{grammar} правило, где левая часть --- имя шаблона, а правая часть --- регулятрное выражение построенное по шаблону.

$$
\infer[]{\langle db, graph, grammar, res \rangle \xrightarrow{\textit{NamedPattern} \ name \ pattern} \langle db, graph, \textit{addRule(grammar, name, pattern)}, res \rangle}{}
$$

Так как теперь запрос может быть не только регулярным, то изменится обработка кинструкции \verb|Query|. Теперь вместо детерминированного конечного автомата строится рекурсивный конечный автомат с учётом накопленной грамматики.

$$
\infer[]{\langle db, graph, grammar, res \rangle \xrightarrow{\textit{Query} \ pattern} \langle db, \textit{buildRSM}(pattern, grammar), grammar, res \rangle}{}
$$

Как следствие, изменяется обработка инструкции \verb|Intersect|: нам необхоимо отслеживать, какого типа графы пересекаются, так как для успешного выполнения операции необходимо, чтобы хотя бы один из них был обыкновенным конечным автоматом. 

$$
\infer[\textit {$g'$ is DFA or $g''$ is DFA}]{\langle db, g, grm, res\rangle \xrightarrow{\textit{Intersect} \ g_1 \ g_2} \langle db, \textit{intersect}(g', g''), grm, res\rangle}{\langle db, g, grm, res\rangle \xrightarrow{g_1} \langle db, g', grm, res \rangle, \ \ \langle db, g, grm, res \rangle \xrightarrow{g_2} \langle db, g'', grm, res \rangle}
$$

Добавим операцию, позволяющую явно указывать стартовые и финальные состояния \verb|SetStartAndFinal|.

{
\scriptsize
$$
\infer[res' \neq \varnothing, res'' \neq \varnothing]{\langle db, g, grm, res\rangle \xrightarrow{\textit{SetStartAndFinal} \ start \ final \ g_1} \langle db, \textit{(g' with Start = $res'$, Final = $res''$)}, grm, res\rangle}{\langle db, g, grm, res\rangle \xrightarrow{g_1} \langle db, g', grm, res \rangle, \ \ \langle db, g', grm, res \rangle \xrightarrow{start} \langle db, g', grm, res' \rangle, \ \ \langle db, g', grm, res \rangle \xrightarrow{final} \langle db, g', grm, res'' \rangle}
$$
}

{
\scriptsize
$$
\infer[res' = \varnothing, res'' \neq \varnothing]{\langle db, g, grm, res\rangle \xrightarrow{\textit{SetStartAndFinal} \ start \ final \ g_1} \langle db, \textit{(g' with Final = $res''$)}, grm, res\rangle}{\langle db, g, grm, res\rangle \xrightarrow{g_1} \langle db, g', grm, res \rangle, \ \ \langle db, g', grm, res \rangle \xrightarrow{start} \langle db, g', grm, res' \rangle, \ \ \langle db, g', grm, res \rangle \xrightarrow{final} \langle db, g', grm, res'' \rangle}
$$
}

{
\scriptsize
$$
\infer[res' \neq \varnothing, res'' = \varnothing]{\langle db, g, grm, res\rangle \xrightarrow{\textit{SetStartAndFinal} \ start \ final \ g_1} \langle db, \textit{(g' with Start = $res$)}, grm, res\rangle}{\langle db, g, grm, res\rangle \xrightarrow{g_1} \langle db, g', grm, res \rangle, \ \ \langle db, g', grm, res \rangle \xrightarrow{start} \langle db, g', grm, res' \rangle, \ \ \langle db, g', grm, res \rangle \xrightarrow{final} \langle db, g', grm, res'' \rangle}
$$
}


{
\scriptsize
$$
\infer[res' = \varnothing, res'' = \varnothing]{\langle db, g, grm, res\rangle \xrightarrow{\textit{SetStartAndFinal} \ start \ final \ g_1} \langle db, g', grm, res\rangle}{\langle db, g, grm, res\rangle \xrightarrow{g_1} \langle db, g', grm, res \rangle, \ \ \langle db, g', grm, res \rangle \xrightarrow{start} \langle db, g', grm, res' \rangle, \ \ \langle db, g', grm, res \rangle \xrightarrow{final} \langle db, g', grm, res'' \rangle}
$$
}

$$
\infer[]{\langle db, graph, grammar, res \rangle \xrightarrow{\textit{Set} \ s} \langle db, g, grammar, s \cap \textit{graph.Vertices} \rangle}{}
$$


$$
\infer[]{\langle db, graph, grammar, res \rangle \xrightarrow{\textit{None}} \langle db, g, grammar, \varnothing \rangle}{}
$$

$$
\infer[]{\langle db, graph, grammar, res \rangle \xrightarrow{\textit{Range} \ x \ y} \langle db, g, grammar, \textit{graph.Vertices} \cap \{ \min(x,y) \ldots \max(x,y)\} \rangle}{}
$$

Расширим возможности обработки результата пересечения. Для этого добавим фильтрацию рёбер. Обработка конструкции \verb|Filter| выглядит следующим образом.

$$
\infer[\textit{type of res' is Set<Edge>}]{\langle db, graph, grammar, res \rangle \xrightarrow{\textit{Filter} \ cond \ src} \langle db, g, grammar, \textit{Set.filter cond res'} \rangle}{\langle db, graph, grammar, res \rangle \xrightarrow{src} \langle db, g, grammar, res' \rangle}
$$



\section{Лекция 12}
 
Парадигмы программирования. Структурное программирование: машины Тьюринга, архитектура фон Неймана, языки-представители.

Начать с автомата (трансдьюсера). Сложение в унарной систее исчисления.

Машина Тьюринга. Машина Поста, нормальные алгорифмы Маркова.

Машина Тьюринга $M=\langle Q,\Gamma ,b,\Sigma ,\delta ,q_{0},F\rangle$:
\begin{itemize}
\item $Q$ --- конечное не пустое множество состояний;
\item $\Gamma$  --- конечное не пустое  множество символов (алфавит);
\item $b \in \Gamma$  специальный пробельный символ;
\item $ \Sigma \subseteq \Gamma \setminus \{b\}$ входной алфавит (то, что можно написать на ленту ``снаружи'' перед стартом);
\item $q_{0}\in Q$ --- стартовое состояние;
\item $F\subseteq Q$ Множество финальных состяний. Говорят, что вход принимается машиной $M$ ечли она оказывается в одном из финальных состояний $F$.
\item $ \delta :(Q\setminus F)\times \Gamma \to Q\times \Gamma \times \{L,R\}$ --- частично определённая функция переходоа. $L$ и $R$ --- команды перемещения головки (влево и вправо соответственно). Если $\delta$  не определена, то машина останавливается.
\end{itemize}

Эмулятор машин Тьюринга: \url{https://turingmachinesimulator.com/}

Архитектура фон Неймана (Принстон). 
\begin{itemize}
	\item Однородность памяти: данные и команды --- одно и то же, хранятся в одной памяти, обрабатываются по общим принципам.
	\item Адресность: память --- это набор занумерованных ячеек.
	\item Программное управление: все действия описываются программой.
\end{itemize}

Узкое место --- канал передачи данных.

Гарвард.
\begin{itemize}
	\item данные и команды разделены. 
\end{itemize}

Проблема: сложнее и дороже. 

Классика: TTA (\url{http://openasip.org/}).




\subsection{Busy beaver}

\begin{itemize}
	\item У машины $n$ состояний + 1 финальное
    \item У машины одна бесконечная в обе стороны лента
    \item Алфавит ленты: $\{0, 1\}$, 0 --- пробельный символ.
    \item Функция перехода получает состояние и текущий символ на ленте, переходит в новое состояние, пишет что-то на ленту, сдвигает головку (налево или направо).
\end{itemize}


Задача: найти терминирующуюся машинутьюринга, записывающую максимальное число 1 на ленту.

Невычислимая функция. 



\subsection{Лямбда-исчисление}

Применение и абстракция. Констант не предполагается.


Сложене в арифметике Пеано.

$\alpha$-эквивалентность --- переименование переменных.

$\beta$-редукция --- вычисление. $(\alpha x.t) a = t[x:=a]$   

$\eta$-преобразование --- заворачивание-разворачивание функций. $\lambda x . f x = f$. Если нет свободных вхождений $x$ в $f$ 

Лямбда-процессоры.

Reduceron: \url{https://www.cs.york.ac.uk/fp/reduceron/}
A Platform for Full-Stack Functional Programming: \url{https://ieeexplore.ieee.org/abstract/document/9180772}
ACQuA: A Parallel Accelerator Architecture for Pure Functional Programs: \url{https://ieeexplore.ieee.org/abstract/document/9155051}

Рекурсивные функции (Гёдель, теория вычислимости).

Кодирование по Чёрчу (Church encoding).

$$
\begin{array}{r|l|l}{\text{Number}}&{\text{Function definition}}&{\text{Lambda expression}}\\\hline 0&0\ f\ x=x&0=\lambda f.\lambda x.x\\1&1\ f\ x=f\ x&1=\lambda f.\lambda x.f\ x\\2&2\ f\ x=f\ (f\ x)&2=\lambda f.\lambda x.f\ (f\ x)\\3&3\ f\ x=f\ (f\ (f\ x))&3=\lambda f.\lambda x.f\ (f\ (f\ x))\\\vdots &\vdots &\vdots \\n&n\ f\ x=f^{n}\ x&n=\lambda f.\lambda x.f^{\circ n}\ x
\end{array}
$$


$$
\operatorname {plus} \equiv \lambda m.\lambda n.\lambda f.\lambda x.m\ f\ (n\ f\ x)
$$

$$
\operatorname {succ} \equiv \lambda n.\lambda f.\lambda x.f\ (n\ f\ x)
$$

$$
 \displaystyle \operatorname {mult} \equiv \lambda m.\lambda n.\lambda f.\lambda x.m\ (n\ f)\ x
$$

Функциональное программирование. Понятие лямбда-исчисления, основные принципы и особенности функционального программирования. Языки представители, Haskell, F\#, Ocaml.


%Программирование в зависимых типах. Изоморфизм Карри-Говарда...
    

\subsection{Парадигмы}
Объектно-ориентированное программирование, основные понятия, инкапсуляция, наследование, полиморфизм. Языки-представители. Пример объектно-ориентированного кода на F\#. 

    













    Логическое программирование, Пролог. 
    SWI-prolog:\url{https://swish.swi-prolog.org/}
    \begin{verbatim}
a(1,2).
a(2,3).
a(3,4).
a(4,1).
b(1,5).
b(5,6).
b(6,1).

reachable(X,Y) :- a(X,Z),b(Z,Y).
reachable(X,Y) :- a(X,Z),b(W,Y),reachable(Z,W).
    \end{verbatim}

функционально-логический (лямбда-пролог, ещё кто-то)

    
%    Рекурсивное программирование, Рефал. 

    Стековое программирование, Форт. 
Forth (Форт): \url{https://www.forth.com/resources/forth-programming-language/}


Визуальное программирование, визуальное моделирование, UML, предметно-ориентированное моделирование.

UML examples: \url{https://www.uml-diagrams.org/index-examples.html}

Trik Studio: \url{https://github.com/trikset/trik-studio}
MetaEdit+: \url{https://www.metacase.com/download/metaedit/moremacosx45.html}





\section{Практика 10}

\begin{enumerate}
    \item \textbf{[5 балла]} Оформление калькулятора или регулярных выражений как отдельного проекта. Создать репозиторий, снабдить всеми необходимыми элементами экосистемы: сборка, тесты, лицензия, readme.
    \item \textbf{[5 баллов]} Создать документацию. Описать цели и задачи проекта, конкретный синтаксис языка, привести примеры.
    \item \textbf{[3 балла]} Создать диаграмму (наиболее подходящего типа), описывающую структуру проекта, выбранного выше. Добавить её в документацию.
\end{enumerate}

\section{Практика 11}

Мини-IDE для языка из предыдущей работы. Проект оформляется в отдельном репозитории по всем правилам. Процесс сдачи задач прежний.

В качестве интерпретатора используется пакет напарника. Все необходимые улучшения интерпретатора, добавление новых функций, исправления оформляются как issues на GitHub в соответсвующес проекте.

\begin{enumerate}
    \item \textbf{[2 балла]} Расширить синтаксис соответствующего языка логическими выражениями (с переменными) и условными операторами.
    \item \textbf{[9 баллов]} Разработать среду разработки для полученного языка. Среда должна предоставлять следующие возможности:
    \begin{itemize}
        \item Редактировать код.
        \item Работать с файлами: создать новый, открыть существующий, сохранить изменения. 
        \item Выводить сообщения о (синтаксических) ошибках.
        \item Запустить программу на исполнение.
        \item Видеть результат исполнения в «консоли».
    \end{itemize}
    \item \textbf{[3 балла]} Расширить IDE возможностью подсветки синтаксиса.
    \item \textbf{[8 баллов]} Расширить IDE возможностью устанавливать точки останова. В момент остановки должна быть возможность просмотреть значения всех ``живых'' переменных.
  
\end{enumerate}
\section{Практика 12}

\subsection{Григорьев С.В.}


Исполнение скриптов запросов.

\href{https://packages.ubuntu.com/search?keywords=antlr4}{ANTLR-пакет для Ubuntu}.

\href{https://github.com/shmatov/antlr4-calculator}{Пример калькулятора}.


Задача на дом.
\begin{enumerate}
  \item Полная поддержка языка запросов. 
  \item Тесты. Очень много тестов. Каждая смысловая конструкция должна бять проверена. 

\end{enumerate}
\section{Практика 11}

\begin{enumerate}
    \item \textbf{[7 баллов]} На основе Mailboxprocessor или Hopac реализовать решение, в котором есть следующие конкурентно выполняющиеся типы задачи: 
    \begin{itemize}
        \item Подгрузка пар матриц из файлов (для генерации использовать генератор из предыдущей работы)
        \item Различные алгоритмы перемножения матриц (для разреженных, для плотных параллельно и последовательно)
        \item Балансировщик, знающий, кому какие матрицы отправлять для обработки.
    \end{itemize}
    Предусмотреть два режима работы: 
    \begin{itemize}
        \item Обработать все матрицы, доступные на входе
        \item Обработать заданное количество пар матриц
    \end{itemize}
    \item \textbf{[14 баллов]} Проанализировать масшатабируемость полученной системы. Какое количество конкурентных задач оптимально для определённой конфигурации системы? Оформить соответствующий отчёт.
\end{enumerate}





\section{САКОД 1: Основы обработки изображений}

Форматы изображений.

Векторный, растровый.

Нас интересует растровый.

И для простоты сразу битмапа (bmp): двумерный массив пикселей, где для каждого хранится цвет.
Естественно, ещё и метаданные вокруг, но они, в основном, про то, как ситать файл, а не про само изображение.
Цвет --- либо RGB, либо градации серого (Grayscale). На пиксель от 1 до 64 бит. В grayscale 16 или 32.

\href{Bitmap in .NET}{https://docs.microsoft.com/en-us/dotnet/api/system.drawing.bitmap?view=dotnet-plat-ext-5.0}


Фильтры.

Собель для поиска границ, размытие по Гауссу, машинное обучение


Реализовать приложение с графическим интерфейсом пользователя, позволяющее открыть папку с изображениями, выбрать изображение, просмотреть его, просмотреть информацию о нём (размер в пикселях, размер в мегабайтах). 

    1. (1 балл) Расширить приложение из предыдущей домашней работы графической компонентой задания матричного фильтра. Необходимо предусмотреть возможность выбора типа фильтра, дефолтных значений, размера фильтра, корректировку весов.
    2. (1 балл) Расширить приложение из предыдущей работы возможностью отображать одновременно два изображения: до и после применения фильтра. Предусмотреть возможность сохранять результат применения фильтра.
    3. (3 балла) Реализовать применение матричных фильтров с использованием GPGPU. Интегрировать с разработанным графическим интерфейсом. Предусмотреть возможность применения нескольких фильтров последовательно.
    4. (6 баллов) Расширить разрабатываемое приложение возможностью потоковой обработки изображений: выбираем папку с изображениями и ко всем применяем заданный фильтр. Результаты применения фильтров сохраняются.
    5. (5 баллов) Подготовить отчёт с анализом производительности и масштабируемости полученного решения.


    Домашки:
    --- гитхаб
    --- ридми
    --- тесты
    --- CI
    --- build-script
    --- ubuntu

\bibliographystyle{abbrv}
\bibliography{Formal_language_course}


\end{document}
